\section{Recursive Least Squares Learning}\label{sec:rec_ls}
% =====================================================

We first consider Recursive Least Squares (RLS) learning, which is a recursive implementation of standard Ordinary Least Squares (OLS). At each time $t$, agents update their estimate $\beta_t$ by minimizing the sum of squared forecast errors:\footnote{We use the following timing: at time $t$ the agent forecasts $r^e_{t+1} = x_t \beta_t$ which then implies a realization for $P_t$ and subsequently $r_{t+1}$. Therefore, to avoid simultaneity issues, we assume that when computing the current estimate, the last available log-return observation is $r_{t-1}$.}
\begin{equation}
    \beta_t = \left(\sum^{t-1}_{i=1} x_{i-1}^2 \right)^{-1} \sum^{t-1}_{i=1} x_{i-1} r_i.
\end{equation}

This estimator admits a recursive formulation. Define $M_{t-1} = \frac{1}{t-1} \sum^{t-1}_{i=1} x_{i-1}^2$. Standard algebra, which we defer to Appendix (\ref{app:RLS}) yields the recursions
\begin{equation}
    M_t = M_{t-1} + \frac{1}{t} \left(x_{t-1}^2 - M_{t-1} \right),
\end{equation}
\begin{equation}
    \beta_t = \beta_{t-1} + \left((t-1) M_{t-1} \right)^{-1} \left( x_{t-2} r_{t-1} - x_{t-2}^2 \beta_{t-1} \right).
\end{equation}

Under RLS learning, actual returns evolve according to
\begin{equation}\label{eq:alm_RLS}
    r_{t-1} = \eta_{t-1}
    + g_{\beta_{t-2}}(x_{t-2}) - g_{\beta_{t-1}}(x_{t-1}),
\end{equation}
where $g_\beta(\cdot)$ is defined in \eqref{eq:g_beta_def}.

A natural question is whether this learning process converges to the rational expectations equilibrium. 
Under FIRE, returns in our model are unpredictable and $\beta = 0$. 
However, convergence to the REE requires not only that agents' perceived relationship matches the actual relationship, but that the entire distributions implied by the PLM and ALM coincide. 
When the ALM is nonlinear and the PLM is linear, such convergence is generally not possible. 
The learning literature has instead focused on weaker notions of equilibrium that require consistency only on specific moments rather than on entire distributions \citep{Hommes2014}. 
Following this approach, we define a Cross-Moment Consistent Equilibrium (CMCE) as an equilibrium in which the population regression coefficient in agents' perceived linear model matches the projection of actual returns onto the predictor.


\begin{definition}[Cross-Moment Consistent Equilibrium]\label{def:cmce}
A pair $(\mu,\beta)$ is a \emph{Cross-Moment Consistent Equilibrium (CMCE)} if:
\begin{enumerate}
    \item $\mu$ is a non-degenerate invariant probability measure for the 
    stochastic process
    \[
        r_{t-1}
        = \eta_{t-1}
        + g_{\beta}(x_{t-2})
        - g_{\beta}(x_{t-1}),
    \]
    where $x_t \stackrel{iid}{\sim} \mathcal{N}(0, \sigma^2_x)$.
    
    \item The parameter $\beta$ is the population OLS coefficient of $r_{t-1}$ 
    on $x_{t-2}$ under the invariant measure $\mu$, that is,
    \[
        \beta = \frac{\operatorname{Cov}_\mu(x_{t-2} r_{t-1})}{\operatorname{Var}_\mu(x_{t-2}^2)} 
    \]
\end{enumerate}
\end{definition}

While a rational expectations equilibrium requires that agents' perceived law of motion coincides with the actual law of motion, a CMCE requires only that the population regression coefficient in agents' perceived linear model matches the projection of actual returns onto the predictor $x_{t-2}$. 
This weaker notion of consistency is appropriate when agents use misspecified linear models to forecast returns generated by a nonlinear process. Despite the misspecification, there is a statistical self-consistency between the empirical relationship agents perceive and the one implied by the ALM.

\begin{proposition}\label{prop:RLS}
    Under RLS learning there exists a unique CMCE given by $\beta = 0$. The CMCE is locally stable under learning for $\kappa \in (0,1)$.
\end{proposition}

\textit{Proof.} In Appendix \ref{app:proof_RLS}.

Notice that $\kappa = \delta a^{-\ell} \phi$ is strictly positive for all admissible parameter values. The restriction $\kappa < 1$ imposes an upper bound on the combination of the discount factor, risk aversion, and dividend growth, and is required for the REE to be well-defined. This condition holds for standard parameter calibrations.

Proposition \ref{prop:RLS} shows that with infinite memory, learning eliminates perceived predictability: agents eventually accumulate enough data to recognize that returns are fundamentally unpredictable.
However, this convergence relies critically on the assumption that all historical observations receive equal weight in estimation. 
In practice, investors may discount older data either because they believe the economic environment is changing or because they face cognitive or computational constraints on memory. 
We now examine how limited memory alters the learning outcome.




% =====================================================
